%! Rational Functions and Zeros — Unit 1, Lesson 1.8 — Exit Ticket (Answer Key)
\documentclass[10pt]{article}
\usepackage{precalculus-article}
\usepackage{precalculus-key}   % pulls in -boxes; do NOT also load -boxes
\newcommand{\TallMath}[1]{$\displaystyle #1\rule[-1.4em]{0pt}{3.2em}$}

\begin{document}

\pageheader{Unit 1, Lesson 1.8}{Exit Ticket --- Answer Key}
\namedateperiod

\vspace{0.14in}

\begin{enumerate}[itemsep=16pt]

\item Let $\displaystyle r(x) = \frac{x^2 - x - 6}{x + 4}$.

      \begin{enumerate}[label=(\alph*), itemsep=10pt]
        \item What value(s) of $x$ are \textbf{not} in the domain of $r$?
              Explain briefly.

              \vspace{0.06in}
              $x = $ \ans{$-4$} because \ansline{$x + 4 = 0$ at $x = -4$, making the denominator zero and $r(-4)$ undefined.}

        \item Find all real zeros of $r(x)$.

              \vspace{0.06in}
              Factor the numerator: $x^2 - x - 6 = $ \ans{$(x - 3)(x + 2)$}

              \vspace{0.06in}
              Numerator zeros: $x = $ \ans{$3$ and $x = -2$}

              \vspace{0.06in}
              Check each against the domain restriction:
              \ansline{Neither $x = 3$ nor $x = -2$ equals $-4$, so both are in the domain.}

              \vspace{0.06in}
              Zeros of $r$: $x = $ \ans{$3$ and $x = -2$}
      \end{enumerate}

\item A student says: ``$x = -4$ is a zero of
      $\displaystyle r(x) = \frac{x + 4}{x^2 - 1}$ because the numerator equals zero
      when $x = -4$.''

      Is the student correct? Circle one: \quad \textbf{No}

      \vspace{0.06in}
      \ansline{The student is \textbf{incorrect}. The numerator $x + 4 = 0$ at $x = -4$, but we also need $x = -4$ to be in the domain of $r$. The denominator $x^2 - 1 = (-4)^2 - 1 = 15 \ne 0$, so $x = -4$ \textbf{is} in the domain. Therefore $r(-4) = 0$, and the student's conclusion is actually correct for the wrong reason --- but the reasoning is incomplete without checking the domain.}

\end{enumerate}

\end{document}
